%============================================================
%  TNChatbot — Rapport de Projet de Fin d'Année
%  Institut International de Technologie à Sfax (IIT)
%  Template fidèle au modèle PFA-5.pdf
%============================================================
\documentclass[12pt,a4paper]{report}

%------------------------------------------------------------
% Packages
%------------------------------------------------------------
\usepackage[utf8]{inputenc}
\usepackage[T1]{fontenc}
\usepackage[french]{babel}
\usepackage{geometry}
\usepackage{graphicx}
\usepackage{fancyhdr}
\usepackage{titlesec}
\usepackage{tocloft}
\usepackage[hidelinks]{hyperref}
\usepackage{array}
\usepackage{booktabs}
\usepackage{tabularx}
\usepackage{longtable}
\usepackage{xcolor}
\usepackage{listings}
\usepackage{float}
\usepackage{setspace}
\usepackage{microtype}
\usepackage{caption}
\usepackage{subcaption}
\usepackage{enumitem}
\usepackage{tikz}
\usepackage{amsmath}
\usepackage{mdframed}
\usepackage{minitoc}
\usepackage{lettrine}
\usepackage{lmodern}
\usepackage{textcomp}
\usepackage{multirow}
\usepackage{colortbl}

%------------------------------------------------------------
% Geometry
%------------------------------------------------------------
\geometry{
  a4paper,
  top=2.5cm, bottom=2.5cm,
  left=3cm,  right=2.5cm,
  headheight=15pt
}

%------------------------------------------------------------
% Line spacing
%------------------------------------------------------------
\setstretch{1.15}

%------------------------------------------------------------
% Colors
%------------------------------------------------------------
\definecolor{chapterbg}{RGB}{0,0,0}
\definecolor{codeframe}{RGB}{200,200,200}
\definecolor{codebg}{RGB}{248,248,248}
\definecolor{keyword}{RGB}{0,0,180}
\definecolor{comment}{RGB}{100,100,100}
\definecolor{string}{RGB}{180,0,0}
\definecolor{tableheader}{RGB}{230,230,230}

%------------------------------------------------------------
% Chapter style — boîte noire + numéro + ligne horizontale
% Fidèle au template IIT
%------------------------------------------------------------
\titleformat{\chapter}[display]
  {\normalfont}
  {%
    \begin{tikzpicture}[baseline=(num.base)]
      \node[fill=black, minimum width=1.4cm, minimum height=1.4cm,
            text=white, font=\Huge\bfseries, inner sep=0pt,
            anchor=south west] (num) at (0,0) {\thechapter};
      \node[font=\small\scshape, anchor=south west, text=black]
            at (-0.05cm,1.45cm) {Chapitre};
      \draw[line width=0.8pt] (1.55cm,0.7cm) -- (\textwidth,0.7cm);
    \end{tikzpicture}%
  }
  {0pt}
  {\vspace{6pt}\Large\bfseries}
  [\vspace{4pt}]

\titleformat{name=\chapter,numberless}[display]
  {\normalfont\Large\bfseries\centering}
  {}
  {0pt}
  {}

\titlespacing*{\chapter}{0pt}{-10pt}{20pt}

%------------------------------------------------------------
% Section / subsection style
%------------------------------------------------------------
\titleformat{\section}{\large\bfseries}{\thesection}{1em}{}
\titleformat{\subsection}{\normalsize\bfseries}{\thesubsection}{1em}{}
\titleformat{\subsubsection}{\normalsize\bfseries}{\thesubsubsection}{1em}{}

%------------------------------------------------------------
% Headers / footers
%------------------------------------------------------------
\pagestyle{fancy}
\fancyhf{}
\renewcommand{\headrulewidth}{0.4pt}
\fancyhead[L]{\small\scshape\leftmark}
\fancyhead[R]{\small\thepage}
\fancyfoot{}
\renewcommand{\chaptermark}[1]{\markboth{\chaptername\ \thechapter.\ #1}{}}

\fancypagestyle{plain}{
  \fancyhf{}
  \fancyfoot[C]{\small\thepage}
  \renewcommand{\headrulewidth}{0pt}
}

%------------------------------------------------------------
% Caption style (italic, small, "–" separator)
%------------------------------------------------------------
\captionsetup{
  labelfont=sc,
  textfont=it,
  labelsep=endash,
  font=small
}

%------------------------------------------------------------
% List style (em-dash bullets like in template)
%------------------------------------------------------------
\setlist[itemize]{label=\textemdash, leftmargin=2em, itemsep=2pt}
\setlist[enumerate]{leftmargin=2em, itemsep=2pt}

%------------------------------------------------------------
% Code listing style
%------------------------------------------------------------
\lstset{
  basicstyle=\ttfamily\footnotesize,
  keywordstyle=\color{keyword}\bfseries,
  commentstyle=\color{comment}\itshape,
  stringstyle=\color{string},
  backgroundcolor=\color{codebg},
  frame=single,
  frameround=tttt,
  rulecolor=\color{codeframe},
  breaklines=true,
  breakatwhitespace=true,
  showstringspaces=false,
  tabsize=4,
  captionpos=b,
  numbers=left,
  numberstyle=\tiny\color{comment},
  numbersep=8pt,
  xleftmargin=12pt
}

%------------------------------------------------------------
% minitoc (Plan en début de chapitre)
%------------------------------------------------------------
\mtcsettitle{minitoc}{Plan}
\mtcsetdepth{minitoc}{2}
\setcounter{minitocdepth}{2}

%------------------------------------------------------------
% TOC / LOF / LOT styling (dots like template)
%------------------------------------------------------------
\renewcommand{\cftchapfont}{\bfseries}
\renewcommand{\cftsecfont}{}
\renewcommand{\cftsubsecfont}{}
\renewcommand{\cftdotsep}{1}

%------------------------------------------------------------
% Metadata
%------------------------------------------------------------
\hypersetup{
  pdftitle={TNChatbot — Rapport PFA},
  pdfauthor={Beya Abid},
  pdfsubject={Chatbot intelligent — Tunisie Numérique}
}

%============================================================
\begin{document}
\dominitoc
%============================================================

%------------------------------------------------------------
% PAGE DE TITRE
%------------------------------------------------------------
\begin{titlepage}
  \begin{center}
    % En-tête
    \noindent
    \begin{minipage}[t]{0.30\textwidth}
      \raggedright
      \footnotesize
      \textbf{Ministère de l'Enseignement Supérieur et de\\
      la Recherche Scientifique}\\[4pt]
      \textbf{Institut International de Technologie à Sfax}
    \end{minipage}
    \hfill
    \begin{minipage}[t]{0.25\textwidth}
      \centering
      % Logo IIT (remplacer par le vrai logo si disponible)
      \fbox{\parbox{3cm}{\centering\vspace{8pt}\textbf{IIT}\\
      \small Institut\\International\\Technologie\vspace{8pt}}}
    \end{minipage}
    \hfill
    \begin{minipage}[t]{0.30\textwidth}
      \raggedleft
      \footnotesize
      \textbf{وزارة التعليم العالي والبحث العلمي}\\[4pt]
      \textbf{المدرسة العليا الدولية الخاصة للتكنولوجيا بصفاقس}
    \end{minipage}

    \vspace{0.5cm}
    \noindent\rule{\textwidth}{0.8pt}
    \vspace{1.5cm}

    {\Huge\bfseries Projet de Fin d'Année}

    \vspace{2cm}

    {\large\itshape Réalisé par :}\\[0.5cm]
    {\Large\bfseries Beya Abid}

    \vspace{2cm}
    \noindent\rule{0.6\textwidth}{0.6pt}\\[0.4cm]
    {\LARGE\bfseries TNChatbot : Chatbot intelligent\\[0.3cm]
    pour Tunisie Numérique}\\[0.4cm]
    \noindent\rule{0.6\textwidth}{0.6pt}

    \vspace{1.5cm}

    {\normalsize\itshape Encadré par :}\\[0.3cm]
    {\normalsize\bfseries [Nom de l'encadrant(e)]}

    \vfill

    \fbox{\parbox{0.5\textwidth}{\centering\vspace{4pt}
    DÉPARTEMENT INFORMATIQUE\vspace{4pt}}}

    \vspace{0.8cm}
    {\bfseries Effectué le \today}

    \vspace{0.5cm}
    {\bfseries Année Universitaire 2024/2025}
  \end{center}
\end{titlepage}

%------------------------------------------------------------
% Numérotation romaine pour les pages préliminaires
%------------------------------------------------------------
\pagenumbering{roman}
\setcounter{page}{1}

%------------------------------------------------------------
% REMERCIEMENTS
%------------------------------------------------------------
\chapter*{Remerciements}
\addcontentsline{toc}{chapter}{Remerciements}

Nous souhaitons adresser nos sincères remerciements à l'équipe de \textbf{Tunisie Numérique}
pour la confiance accordée dans le cadre de ce projet et pour la richesse des échanges qui ont
permis d'orienter les choix fonctionnels et techniques.

Nous remercions également notre encadrant(e) académique pour sa disponibilité, ses
conseils avisés et la rigueur avec laquelle il/elle a suivi l'avancement du projet. Sa capacité à
guider notre réflexion dans les moments de doute a été une aide précieuse tout au long de cette
expérience.

Enfin, nous exprimons notre gratitude envers l'Institut International de Technologie
de Sfax pour la qualité de la formation dispensée, ainsi qu'envers les membres du jury qui ont
accepté d'évaluer ce travail.

%------------------------------------------------------------
% RÉSUMÉ
%------------------------------------------------------------
\chapter*{Résumé}
\addcontentsline{toc}{chapter}{Résumé}

Ce rapport présente la conception et le développement de \textbf{TNChatbot}, un assistant
conversationnel intelligent destiné à la plateforme de médias numériques \textbf{Tunisie
Numérique}. L'objectif principal est d'automatiser la qualification commerciale des prospects
en leur proposant un parcours guidé permettant de découvrir les offres publicitaires, d'obtenir
une recommandation selon leur budget, et de laisser leurs coordonnées à l'équipe commerciale.

La solution repose sur une architecture microservices déployée via Docker Compose,
combinant un frontend \textbf{Next.js}, un backend \textbf{FastAPI}, une base de données
\textbf{PostgreSQL}, un moteur de recherche vectorielle \textbf{Qdrant} et un modèle de
langage local via \textbf{Ollama}. Un pipeline \textbf{RAG} (\textit{Retrieval-Augmented
Generation}) hybride — combinant recherche vectorielle et lexicale BM25 — garantit la
pertinence des réponses factuelles.

Le projet a été réalisé en trois sprints selon une méthodologie Kanban : mise en place
de l'infrastructure et du pipeline d'ingestion, amélioration de la qualité de l'IA (classifieur
d'intention hybride, découpage sémantique, RAG conditionnel), puis renforcement de la
sécurité par authentification JWT et mise en place d'une suite de tests de non-régression.

\textbf{Mots-clés :} Chatbot, RAG, Retrieval-Augmented Generation, LLM, FastAPI,
Next.js, Qdrant, BM25, Docker, Intent Classification, NLP.

%------------------------------------------------------------
% TABLE DES MATIÈRES
%------------------------------------------------------------
\tableofcontents
\addcontentsline{toc}{chapter}{Table des matières}

%------------------------------------------------------------
% TABLE DES FIGURES
%------------------------------------------------------------
\listoffigures
\addcontentsline{toc}{chapter}{Table des figures}

%------------------------------------------------------------
% LISTE DES TABLEAUX
%------------------------------------------------------------
\listoftables
\addcontentsline{toc}{chapter}{Liste des tableaux}

%------------------------------------------------------------
% Numérotation arabe pour le corps du document
%------------------------------------------------------------
\pagenumbering{arabic}
\setcounter{page}{1}

%============================================================
% INTRODUCTION GÉNÉRALE
%============================================================
\chapter*{Introduction Générale}
\addcontentsline{toc}{chapter}{Introduction Générale}
\markboth{Introduction Générale}{}

\lettrine[lines=2]{D}{ans} un contexte de transformation digitale accélérée, les plateformes
de médias en ligne font face à un défi croissant : répondre efficacement aux sollicitations
commerciales d'un nombre grandissant de prospects tout en maintenant la qualité du service.
\textbf{Tunisie Numérique}, l'un des portails d'information numérique les plus visités en
Tunisie, n'échappe pas à cette réalité.

Aujourd'hui, le processus de découverte des offres publicitaires — bannières display,
contenus sponsorisés, vidéos, newsletters — repose encore largement sur des échanges manuels
entre les prospects et l'équipe commerciale. Ce modèle génère un volume important de
questions répétitives, ralentit la qualification des leads et monopolise des ressources humaines
qui pourraient être consacrées à des interactions à plus forte valeur ajoutée.

C'est dans ce cadre que s'inscrit le projet \textbf{TNChatbot} : concevoir et déployer un
assistant conversationnel intelligent capable d'informer, d'orienter et de qualifier les prospects
de manière autonome, tout en garantissant la pertinence des réponses grâce à une base de
connaissances structurée.

La solution proposée combine un parcours conversationnel guidé, un moteur de réponse
enrichi par la technologie \textit{Retrieval-Augmented Generation} (RAG), et une logique de
qualification commerciale via des formulaires de capture de leads. L'ensemble est déployable
localement via Docker Compose, ce qui garantit la reproductibilité des environnements et
facilite la mise en production.

Ce rapport est organisé comme suit. Le \textbf{chapitre 1} pose les bases du projet en
analysant le contexte, la problématique, les solutions existantes et la méthodologie adoptée.
Le \textbf{chapitre 2} détaille les choix de conception et l'architecture technique retenue.
Enfin, le \textbf{chapitre 3} présente la réalisation concrète du projet, sprint par sprint, en
mettant en lumière les décisions techniques et les résultats obtenus.

%============================================================
% CHAPITRE 1 — ÉTUDE PRÉALABLE
%============================================================
\chapter{Étude Préalable}
\minitoc
\newpage

%------------------------------------------------------------
\section{Introduction}
%------------------------------------------------------------

Ce chapitre pose les bases de notre projet en présentant le contexte opérationnel de
\textbf{Tunisie Numérique}, en identifiant les défis liés à la qualification commerciale
manuelle, et en analysant les solutions conversationnelles existantes sur le marché. Cette
étude préalable nous permet ensuite de définir précisément les objectifs du projet et de
proposer une solution adaptée, avant d'exposer la méthodologie agile adoptée pour la
conduite du développement.

%------------------------------------------------------------
\section{Cadre général du projet}
%------------------------------------------------------------

\subsection{Présentation de Tunisie Numérique}

\textbf{Tunisie Numérique} est l'une des plateformes d'information digitale les plus influentes
en Tunisie. Avec plusieurs millions de visiteurs uniques par mois, le portail couvre l'actualité
politique, économique, culturelle et technologique du pays. Sa régie publicitaire propose
aux annonceurs, agences et entreprises une gamme d'offres variées : bannières
\textit{display}, contenus sponsorisés, vidéos pré-roll, newsletters et solutions premium.

Dans cette logique de croissance commerciale, Tunisie Numérique cherche à améliorer
l'expérience de ses prospects et à optimiser la productivité de son équipe commerciale. Le
chatbot TNChatbot s'inscrit dans cette démarche de transformation digitale.

\subsection{Problématique et motivations}

Le processus de découverte des offres publicitaires de Tunisie Numérique repose
actuellement sur des échanges manuels entre les prospects et l'équipe commerciale. Les
principaux défis identifiés sont les suivants :

\begin{itemize}
  \item \textbf{Volume élevé de questions répétitives} : les prospects posent systématiquement
  les mêmes questions sur les formats disponibles, les audiences et les tarifs.
  \item \textbf{Difficulté à choisir rapidement une solution adaptée} : sans outil d'orientation,
  les prospects peinent à identifier l'offre correspondant à leur budget et leurs objectifs.
  \item \textbf{Qualification insuffisante des leads entrants} : les demandes de contact arrivent
  sans contexte suffisant pour que l'équipe commerciale puisse les prioriser.
  \item \textbf{Absence d'un canal automatisé} : aucun outil ne permet d'orienter les prospects
  en dehors des heures ouvrables.
\end{itemize}

La problématique centrale est donc la suivante : \textit{comment proposer un assistant
conversationnel capable d'informer, d'orienter et de qualifier efficacement les prospects,
tout en garantissant la pertinence des réponses et la valeur commerciale des leads collectés ?}

\subsection{Étude de l'existant}

Plusieurs catégories de solutions conversationnelles existent sur le marché. Le tableau
suivant présente une comparaison selon les critères les plus pertinents pour notre besoin.

\begin{table}[H]
\centering
\caption{Comparaison des solutions conversationnelles existantes}
\small
\begin{tabularx}{\textwidth}{|l|X|X|X|X|}
\hline
\rowcolor{tableheader}
\textbf{Critère} & \textbf{Chatbot à règles} & \textbf{Chatbot FAQ} & \textbf{LLM générique} & \textbf{RAG + LLM} \\
\hline
Pertinence métier & Faible & Moyenne & Moyenne & \textbf{Élevée} \\
\hline
Personnalisation & Faible & Faible & Moyenne & \textbf{Élevée} \\
\hline
Fiabilité factuelle & Moyenne & Élevée & Faible & \textbf{Élevée} \\
\hline
Collecte de leads & Non & Non & Non & \textbf{Oui} \\
\hline
Déploiement local & Oui & Oui & Non & \textbf{Oui} \\
\hline
Coût exploitation & Faible & Faible & Élevé & \textbf{Maîtrisé} \\
\hline
\end{tabularx}
\label{tab:existant}
\end{table}

Les solutions à base de règles ou de FAQ sont limitées par leur incapacité à traiter des
questions non prévues. Les LLM génériques (GPT-4, Claude) offrent de meilleures capacités
de compréhension mais génèrent des hallucinations sur les données métier spécifiques. La
combinaison \textbf{RAG + LLM local} représente la solution la plus adaptée : elle ancre les
réponses dans une base documentaire vérifiée tout en permettant une interaction naturelle.

\subsection{Solution proposée}

Face aux limites des approches existantes, nous proposons \textbf{TNChatbot}, une
solution hybride articulée autour de quatre composantes :

\begin{enumerate}
  \item \textbf{Parcours conversationnel guidé} : une machine à états orchestre le dialogue
  selon les étapes clés (accueil, découverte des offres, recommandation budget, capture
  de lead).
  \item \textbf{Moteur RAG hybride} : les questions factuelles sont traitées via une recherche
  combinant similarité vectorielle (\textit{embeddings}) et recherche lexicale BM25, permettant
  une récupération précise des passages pertinents de la base documentaire.
  \item \textbf{Modèle de langage local} : Ollama héberge un LLM local (compatible
  API Ollama) qui génère des réponses naturelles enrichies par le contexte RAG, sans
  dépendance à des services cloud tiers.
  \item \textbf{Module de qualification commerciale} : des formulaires structurés collectent
  les informations de contact et les transmettent à l'équipe commerciale sous forme de
  leads exploitables.
\end{enumerate}

%------------------------------------------------------------
\section{Méthodologie de travail}
%------------------------------------------------------------

\subsection{Méthode Kanban}

Pour assurer une gestion efficace du projet dans un cadre PFE avec des contraintes de
temps imposées, nous avons adopté une méthodologie \textbf{Kanban agile}. Cette approche
présente plusieurs avantages adaptés à notre contexte :

\begin{itemize}
  \item \textbf{Visualisation du flux de travail} : un tableau Kanban divise les tâches en
  colonnes « À faire », « En cours », « En revue » et « Terminé », offrant une visibilité
  permanente sur l'avancement.
  \item \textbf{Limitation du travail en cours (WIP)} : afin d'éviter la surcharge et de
  maintenir une concentration sur les tâches prioritaires.
  \item \textbf{Amélioration continue} : des points réguliers permettent d'ajuster les
  priorités en fonction des retours de l'équipe Tunisie Numérique.
  \item \textbf{Flexibilité} : la méthode s'adapte facilement aux changements de périmètre
  sans perturber le flux existant.
\end{itemize}

\subsection{Planification des sprints}

Le projet a été découpé en \textbf{trois sprints} selon une approche incrémentale.
Chaque sprint livrait un ensemble de fonctionnalités opérationnelles.

\begin{longtable}{|p{2.2cm}|p{7cm}|p{2.8cm}|}
\caption{Planification des sprints du projet TNChatbot} \\
\hline
\rowcolor{tableheader}
\textbf{Sprint} & \textbf{Objectifs et tâches} & \textbf{Durée} \\
\hline
\endfirsthead
\multicolumn{3}{c}{\tablename\ \thetable\ — Suite} \\
\hline
\rowcolor{tableheader}
\textbf{Sprint} & \textbf{Objectifs et tâches} & \textbf{Durée} \\
\hline
\endhead
\hline
\endfoot
Sprint 0
(Cadrage) &
— Analyse du cahier des charges.\newline
— Définition du périmètre fonctionnel.\newline
— Choix des technologies.\newline
— Mise en place de l'environnement Docker. &
Semaines 1--2 \\
\hline
Sprint 1
(Infrastructure) &
— Configuration Docker Compose (6 conteneurs).\newline
— Pipeline d'ingestion KB (chunking, embeddings, Qdrant).\newline
— API conversationnelle (session, message, stream).\newline
— Machine à états conversationnelle.\newline
— Interface admin (discussions, leads, KB). &
Semaines 3--6 \\
\hline
Sprint 2
(Qualité IA) &
— Classifieur d'intention hybride (règles + regex + score).\newline
— Découpage sémantique par paragraphes.\newline
— RAG conditionnel (\texttt{should\_trigger\_rag}).\newline
— Ingestion idempotente par \texttt{source\_uri}. &
Semaines 7--9 \\
\hline
Sprint 3
(Sécurité \& Qualité) &
— Authentification JWT HS256 (stdlib Python).\newline
— Audit log des actions sensibles.\newline
— Suite de tests non-régression (24 cas de test).\newline
— Mise à jour du frontend pour Bearer token. &
Semaines 10--12 \\
\hline
\label{tab:sprints}
\end{longtable}

%------------------------------------------------------------
\section{Conclusion}
%------------------------------------------------------------

Ce chapitre a posé les fondations du projet en analysant le contexte opérationnel de
Tunisie Numérique, les limites des solutions conversationnelles existantes, et en proposant
une approche RAG + LLM local comme réponse adaptée. La méthodologie Kanban, articulée
en trois sprints progressifs, fournit un cadre structuré pour la réalisation du projet.
Le chapitre suivant détaille les choix de conception et l'architecture technique retenue.

%============================================================
% CHAPITRE 2 — CONCEPTION ET ARCHITECTURE
%============================================================
\chapter{Conception et Architecture}
\minitoc
\newpage

%------------------------------------------------------------
\section{Introduction}
%------------------------------------------------------------

Ce chapitre présente l'architecture générale de TNChatbot, les choix technologiques
justifiés, le modèle de données, la machine à états conversationnelle et le pipeline RAG.
Ces décisions de conception constituent le socle sur lequel s'appuie l'ensemble de la
réalisation.

%------------------------------------------------------------
\section{Architecture générale du système}
%------------------------------------------------------------

TNChatbot repose sur une architecture \textbf{microservices} déployée via
\textbf{Docker Compose}. Six conteneurs coopèrent pour former le système complet.

\begin{figure}[H]
\centering
\begin{tikzpicture}[
  box/.style={draw, rounded corners=4pt, minimum width=2.8cm,
              minimum height=1cm, align=center, font=\small},
  arrow/.style={->, >=stealth, thick}
]
  % User
  \node[box, fill=gray!10] (user) at (0,6) {Utilisateur\\(Navigateur)};

  % Docker host
  \draw[dashed, rounded corners=6pt]
    (-4.5,-1) rectangle (9.5,5.5);
  \node[font=\footnotesize\itshape] at (2.5,5.2) {Docker Host (docker-compose)};

  % Frontend
  \node[box, fill=blue!10] (front) at (-2,3.5) {frontend\\Next.js\\\small :19080};
  \node[box, fill=blue!8] (proxy) at (-2,2) {API Proxy\\route.ts};

  % Backend
  \node[box, fill=green!10] (back) at (2,3.5) {backend\\FastAPI\\\small :8000};
  \node[box, fill=green!8] (orch) at (0.5,1.5) {Orchestrateur\\state\_machine};
  \node[box, fill=green!8] (rag) at (2.5,0.5) {Module RAG\\retrieve.py};
  \node[box, fill=green!8] (leads) at (4.5,1.5) {Admin/\\Leads};
  \node[box, fill=green!8] (llm_c) at (2.5,2.5) {LLM Client};

  % Datastores
  \node[box, fill=orange!10] (pg) at (-1,0) {postgres\\\small :5432};
  \node[box, fill=purple!10] (qd) at (3,-0.7) {qdrant\\\small :6333};
  \node[box, fill=red!10] (ollama) at (6.5,1) {ollama\\\small :11434};
  \node[box, fill=red!8] (init) at (6.5,-0.5) {ollama-init};

  % Arrows
  \draw[arrow] (user) -- (front);
  \draw[arrow] (front) -- (proxy);
  \draw[arrow] (proxy) -- (back) node[midway, above, font=\tiny] {HTTP};
  \draw[arrow] (back) -- (orch);
  \draw[arrow] (back) -- (leads);
  \draw[arrow] (orch) -- (rag);
  \draw[arrow] (back) -- (llm_c);
  \draw[arrow] (back) -- (pg) node[midway, left, font=\tiny] {SQL};
  \draw[arrow] (rag) -- (qd) node[midway, right, font=\tiny] {vecteurs};
  \draw[arrow] (llm_c) -- (ollama) node[midway, right, font=\tiny] {chat};
  \draw[arrow] (init) -- (ollama) node[midway, right, font=\tiny] {pull};
\end{tikzpicture}
\caption{Architecture générale de TNChatbot — vue conteneurs}
\label{fig:archi}
\end{figure}

Le flux est le suivant : l'utilisateur interagit avec l'interface Next.js, qui transmet les
requêtes via un proxy HTTP au backend FastAPI. Ce dernier orchestre la logique
conversationnelle via une machine à états, fait appel au module RAG si nécessaire,
interagit avec PostgreSQL et Qdrant, puis délègue la génération de réponse à Ollama.

%------------------------------------------------------------
\section{Choix technologiques}
%------------------------------------------------------------

\begin{table}[H]
\centering
\caption{Comparaison et justification des choix technologiques}
\small
\begin{tabularx}{\textwidth}{|l|X|X|l|}
\hline
\rowcolor{tableheader}
\textbf{Composant} & \textbf{Technologie retenue} & \textbf{Alternatives étudiées} & \textbf{Justification} \\
\hline
Frontend & Next.js 14 (React) & Vue.js, Angular & SSR, streaming SSE natif \\
\hline
Backend & FastAPI (Python) & Django, Flask, Express & Async natif, typing, OpenAPI \\
\hline
Base SQL & PostgreSQL & MySQL, SQLite & Robustesse, JSON, BM25 \\
\hline
Vectoriel & Qdrant & Pinecone, Weaviate, FAISS & Local, REST API, payload riche \\
\hline
LLM & Ollama & OpenAI API, Mistral & Local, pas de coût API \\
\hline
Conteneurs & Docker Compose & Kubernetes & Simplicité PFE, reproductible \\
\hline
\end{tabularx}
\label{tab:techno}
\end{table}

\subsection{Backend — FastAPI}

FastAPI a été retenu pour sa capacité à gérer nativement les requêtes asynchrones et
le \textit{streaming} SSE (\textit{Server-Sent Events}), indispensables pour l'affichage
progressif des réponses du chatbot. Sa génération automatique de documentation OpenAPI
facilite la communication entre frontend et backend.

\subsection{Base vectorielle — Qdrant}

Qdrant stocke les vecteurs d'embedding associés aux chunks documentaires. Sa capacité
à filtrer les résultats par métadonnées (\texttt{intent}, \texttt{section}) permet
d'affiner le retrieval RAG selon le contexte conversationnel, sans requête secondaire.

\subsection{LLM local — Ollama}

Ollama expose une API compatible OpenAI pour les inférences de chat et la génération
d'embeddings. L'hébergement local garantit la confidentialité des données prospects et
élimine la dépendance à des services cloud tiers, ce qui est essentiel dans un contexte
de données sensibles.

%------------------------------------------------------------
\section{Modèle de données}
%------------------------------------------------------------

Le schéma PostgreSQL organise les données en six tables principales :

\begin{table}[H]
\centering
\caption{Description des tables principales de la base de données}
\small
\begin{tabularx}{\textwidth}{|l|X|}
\hline
\rowcolor{tableheader}
\textbf{Table} & \textbf{Rôle} \\
\hline
\texttt{chat\_sessions} & Une ligne par session utilisateur (id, step, created\_at) \\
\hline
\texttt{chat\_messages} & Historique des messages (role, content, step, session\_id) \\
\hline
\texttt{leads} & Données de contact collectées (nom, société, email, téléphone) \\
\hline
\texttt{kb\_documents} & Métadonnées des documents ingérés (source\_uri, title, status) \\
\hline
\texttt{kb\_chunks} & Chunks et leurs embeddings JSON, liés à kb\_documents \\
\hline
\texttt{kb\_ingestion\_runs} & Traçabilité des runs d'ingestion (status, stats JSON) \\
\hline
\texttt{admin\_config} & Configuration métier clé/valeur (audience, offres, email) \\
\hline
\end{tabularx}
\label{tab:schema}
\end{table}

%------------------------------------------------------------
\section{Machine à états conversationnelle}
%------------------------------------------------------------

L'orchestration du dialogue repose sur une \textbf{machine à états finis} implémentée
dans \texttt{state\_machine.py}. Chaque état correspond à une étape du parcours
utilisateur.

\begin{figure}[H]
\centering
\begin{tikzpicture}[
  state/.style={draw, rounded corners=4pt, minimum width=2.8cm,
                minimum height=0.8cm, align=center, font=\small},
  arrow/.style={->, >=stealth, thick, font=\tiny}
]
  \node[state, fill=gray!15] (welcome) at (0,6) {WELCOME\_SCOPE};
  \node[state] (menu) at (0,4.5) {MAIN\_MENU};
  \node[state, fill=blue!8] (audience) at (-4,3) {AUDIENCE};
  \node[state, fill=blue!8] (solutions) at (0,3) {SOLUTIONS\_MENU};
  \node[state, fill=blue!8] (budget) at (4,3) {BUDGET\_CLIENT\_TYPE};
  \node[state, fill=green!8] (display) at (-3,1.5) {SOLUTION\_DISPLAY};
  \node[state, fill=green!8] (video) at (0,1.5) {SOLUTION\_VIDEO};
  \node[state, fill=orange!8] (form) at (4,1.5) {FORM\_STANDARD\_*};
  \node[state, fill=red!8] (oos) at (-4,4.5) {OUT\_OF\_SCOPE\_READER};

  \draw[arrow] (welcome) -- (menu);
  \draw[arrow] (menu) -- node[left]{audience} (audience);
  \draw[arrow] (menu) -- node[right,near start]{solutions} (solutions);
  \draw[arrow] (menu) -- node[right]{budget} (budget);
  \draw[arrow] (solutions) -- (display);
  \draw[arrow] (solutions) -- (video);
  \draw[arrow] (budget) -- (form);
  \draw[arrow, dashed] (welcome) -- node[above]{\scriptsize lecteur} (oos);
\end{tikzpicture}
\caption{Machine à états conversationnelle (états principaux)}
\label{fig:fsm}
\end{figure}

Les étapes de type \textbf{wizard} (formulaires) et \textbf{statiques} (menus) sont
traitées avant l'appel au pipeline RAG, évitant tout appel d'embedding inutile.

%------------------------------------------------------------
\section{Pipeline RAG}
%------------------------------------------------------------

Le pipeline RAG (\textit{Retrieval-Augmented Generation}) se déroule en huit étapes
lorsqu'une question factuelle est détectée.

\begin{figure}[H]
\centering
\begin{tikzpicture}[
  step/.style={draw, fill=blue!8, rounded corners=4pt,
               minimum width=3cm, minimum height=0.7cm, align=center, font=\small},
  arrow/.style={->, >=stealth, thick}
]
  \node[step] (q) at (0,8) {Question utilisateur};
  \node[step] (intent) at (0,7) {Détection d'intention};
  \node[step] (trigger) at (0,6) {should\_trigger\_rag()};
  \node[step, fill=orange!8] (embed) at (0,5) {Embedding requête};
  \node[step] (vector) at (-3,4) {Recherche Qdrant\\(vectorielle)};
  \node[step] (bm25) at (3,4) {Recherche BM25\\(lexicale)};
  \node[step] (rrf) at (0,3) {Fusion RRF + Reranking};
  \node[step, fill=green!8] (ctx) at (0,2) {Construction contexte RAG};
  \node[step, fill=green!8] (llm) at (0,1) {Génération LLM};
  \node[step] (resp) at (0,0) {Réponse + métadonnées RAG};

  \foreach \from/\to in {q/intent, intent/trigger, trigger/embed,
                          embed/vector, embed/bm25,
                          vector/rrf, bm25/rrf,
                          rrf/ctx, ctx/llm, llm/resp}
    \draw[arrow] (\from) -- (\to);
\end{tikzpicture}
\caption{Pipeline RAG hybride de TNChatbot}
\label{fig:rag}
\end{figure}

\paragraph{Ingestion.}
Les fichiers source (\texttt{.txt}, \texttt{.md}) du dossier \texttt{kb\_sources/} sont
découpés en chunks sémantiques (par paragraphes), vectorisés via Ollama embeddings, puis
stockés dans Qdrant (vecteurs + payload métadonnées) et PostgreSQL (métadonnées + contenu).

\paragraph{Retrieval hybride.}
La recherche combine similarité cosine dans Qdrant et BM25 sur \texttt{kb\_chunks}, les
résultats étant fusionnés via \textit{Reciprocal Rank Fusion} (RRF) puis rerankés
lexicalement (ou via Cohere si la clé API est disponible).

%------------------------------------------------------------
\section{Conclusion}
%------------------------------------------------------------

L'architecture modulaire adoptée — microservices conteneurisés, séparation stricte des
responsabilités, pipeline RAG hybride — garantit à la fois la maintenabilité du système et
sa capacité à évoluer. Le chapitre suivant présente la réalisation concrète, sprint par sprint.

%============================================================
% CHAPITRE 3 — RÉALISATION
%============================================================
\chapter{Réalisation}
\minitoc
\newpage

%------------------------------------------------------------
\section{Introduction}
%------------------------------------------------------------

Ce chapitre présente la réalisation technique du projet, organisée en trois sprints
successifs. Pour chaque sprint, nous détaillons les objectifs, les décisions techniques clés,
les extraits de code les plus significatifs et les résultats obtenus. La dernière section
présente les interfaces utilisateur et administrateur développées.

%------------------------------------------------------------
\section{Sprint 1 — Infrastructure et pipeline d'ingestion}
%------------------------------------------------------------

\subsection{Configuration Docker Compose}

La stack est définie dans \texttt{infra/docker-compose.yml} et comprend six services.
Chaque service est configuré avec des variables d'environnement permettant une
personnalisation sans modification du code source.

\begin{lstlisting}[language=bash, caption={Démarrage de la stack Docker Compose}, label={lst:docker}]
# Démarrage complet de la stack
docker compose -f infra/docker-compose.yml up -d

# Vérification de l'état des conteneurs
docker compose -f infra/docker-compose.yml ps

# Rebuild après modification du code backend
docker compose -f infra/docker-compose.yml build backend
docker compose -f infra/docker-compose.yml up -d backend
\end{lstlisting}

\subsection{Pipeline d'ingestion de la base de connaissances}

Le pipeline d'ingestion (\texttt{ingest.py}) transforme les fichiers sources en vecteurs
exploitables par le moteur RAG. Le processus se déroule en cinq étapes :

\begin{enumerate}
  \item \textbf{Découverte des fichiers} : parcours du dossier \texttt{kb\_sources/}
  (filtres : \texttt{.md}, \texttt{.txt}, hors fichiers de training/synthétiques).
  \item \textbf{Découpage sémantique} : chaque fichier est découpé en chunks respectant
  les frontières de paragraphes (double saut de ligne).
  \item \textbf{Génération des embeddings} : appel à l'API Ollama par batches avec retry.
  \item \textbf{Persistance duale} : métadonnées en PostgreSQL, vecteurs dans Qdrant.
  \item \textbf{Traçabilité} : chaque run d'ingestion est enregistré dans
  \texttt{kb\_ingestion\_runs}.
\end{enumerate}

\subsection{API conversationnelle}

Le backend expose quatre endpoints principaux :

\begin{table}[H]
\centering
\caption{Endpoints de l'API conversationnelle}
\small
\begin{tabularx}{\textwidth}{|l|l|X|}
\hline
\rowcolor{tableheader}
\textbf{Méthode} & \textbf{Route} & \textbf{Rôle} \\
\hline
POST & \texttt{/api/chat/session} & Crée une session, initialise l'étape \\
\hline
POST & \texttt{/api/chat/message} & Traitement synchrone complet (RAG + LLM) \\
\hline
POST & \texttt{/api/chat/stream} & Même logique en streaming SSE \\
\hline
GET & \texttt{/health} & Healthcheck pour discovery admin/frontend \\
\hline
\end{tabularx}
\label{tab:api}
\end{table}

%------------------------------------------------------------
\section{Sprint 2 — Qualité de l'intelligence artificielle}
%------------------------------------------------------------

Le Sprint 2 a ciblé trois dettes techniques majeures identifiées dans le Sprint 1 :
la fragilité du classifieur d'intention, la mauvaise qualité du découpage documentaire,
et le déclenchement systématique du pipeline RAG.

\subsection{Classifieur d'intention hybride}

\subsubsection{Problème initial}

Le classifieur original utilisait un dictionnaire de mots-clés bruts, sensible à la casse
et aux accents, et retournait le premier intent correspondant (ordre arbitraire).
Ce comportement générait des erreurs dès qu'un utilisateur employait un synonyme
ou une graphie alternative (« bannière » au lieu de « banniere »).

\subsubsection{Solution implémentée}

Le nouveau classifieur repose sur la structure \texttt{INTENT\_RULES}, un tableau de
quadruplets \texttt{(intent, keywords, patterns\_regex, weight)}.

\begin{lstlisting}[language=Python, caption={Structure INTENT\_RULES du classifieur hybride},
                   label={lst:intent_rules}]
INTENT_RULES: List[Tuple[str, List[str], List[str], float]] = [
    ("display", [
        "display", "banniere", "bannieres", "banner",
        "leaderboard", "mpu", "encart", "visuel",
    ], [r"banniere.*pub", r"format.*display"], 1.0),
    ("video", [
        "video", "preroll", "instream", "outstream",
        "live", "youtube", "stream",
    ], [r"format.*video", r"pub.*video"], 1.0),
    # ... autres intents
]

def classify_intent(user_message: str) -> Optional[str]:
    normalized = _normalize_text(user_message)  # accents, casse, ponctuation
    best_intent, best_score = None, 0.0
    for intent, keywords, patterns, weight in INTENT_RULES:
        hits = sum(1 for kw in keywords if kw in normalized)
        regex = sum(1 for p in patterns if re.search(p, normalized))
        score = (hits + regex * 2) * weight
        if score > best_score:
            best_score, best_intent = score, intent
    return best_intent
\end{lstlisting}

\paragraph{Normalisation.} La fonction \texttt{\_normalize\_text()} décompose le texte
en NFD (Unicode), supprime les diacritiques et remplace la ponctuation par des espaces.
Ainsi, « Bannières publicitaires » devient « bannieres publicitaires » avant la recherche.

\paragraph{Score pondéré.} Le score est calculé comme
\(\text{score} = (\text{hits\_keywords} + 2 \times \text{hits\_regex}) \times \text{weight}\).
Les patterns regex ont un bonus $\times 2$ pour récompenser les correspondances
contextuelles plus spécifiques. L'intent avec le meilleur score global gagne.

\subsection{Découpage sémantique par paragraphes}

\subsubsection{Problème initial}

La fonction \texttt{chunk\_text()} découpait le texte mot à mot, ignorant complètement
les frontières sémantiques. Un chunk pouvait ainsi contenir la fin d'une section et le
début d'une autre, dégradant la cohérence contextuelle et la pertinence du retrieval RAG.

\subsubsection{Solution implémentée}

\begin{lstlisting}[language=Python, caption={Découpage sémantique par paragraphes},
                   label={lst:chunking}]
def chunk_text(text: str, max_tokens: int, overlap: int) -> List[str]:
    # 1. Découper par frontières naturelles (double newline)
    paragraphs = [p.strip() for p in re.split(r"\n\s*\n", text) if p.strip()]
    chunks, current_words = [], []

    for para in paragraphs:
        para_words = para.split()
        # Paragraphe seul trop long -> découper par mots avec overlap
        if len(para_words) > max_tokens:
            if current_words:
                chunks.append(" ".join(current_words))
                current_words = []
            # Découpage par fenêtre glissante
            start = 0
            while start < len(para_words):
                end = min(start + max_tokens, len(para_words))
                chunks.append(" ".join(para_words[start:end]))
                if end == len(para_words): break
                start = max(0, end - overlap)
        # Ajout dépasse la limite -> flush + overlap
        elif current_words and len(current_words) + len(para_words) > max_tokens:
            chunks.append(" ".join(current_words))
            current_words = (current_words[-overlap:] if overlap else []) + para_words
        else:
            current_words.extend(para_words)

    if current_words:
        chunks.append(" ".join(current_words))
    return [c for c in chunks if c.strip()]
\end{lstlisting}

Le nouveau découpage respecte les unités sémantiques naturelles du texte (paragraphes
Markdown, sections), ce qui améliore significativement la cohérence des chunks indexés.

\subsection{RAG conditionnel}

\subsubsection{Problème initial}

La fonction \texttt{should\_trigger\_rag()} retournait systématiquement \texttt{True},
déclenchant l'appel à l'API d'embeddings et la recherche Qdrant pour \textit{chaque}
message, y compris les simples clics de boutons de navigation.

\subsubsection{Solution implémentée}

\begin{lstlisting}[language=Python, caption={Déclenchement conditionnel du RAG},
                   label={lst:rag_trigger}]
STEPS_NO_RAG = {
    "FORM_STANDARD_EMAIL", "FORM_STANDARD_PHONE", "FORM_STANDARD_DONE",
    "BUDGET_CLIENT_TYPE", "BUDGET_OBJECTIVE", "WELCOME_SCOPE",
    # ... toutes les étapes wizard/formulaire
}

INTENTS_TRIGGERING_RAG = {
    "audience", "display", "video", "content",
    "newsletter_audio", "immoneuf", "premium", "mag",
}

def should_trigger_rag(intent: Optional[str], user_message: str,
                        step: Optional[str] = None) -> bool:
    # Règle 1 : jamais sur étapes wizard/formulaire
    if step and step.upper() in STEPS_NO_RAG:
        return False
    # Règle 2 : toujours si intent documentaire reconnu
    if intent and intent in INTENTS_TRIGGERING_RAG:
        return True
    # Règle 3 : si question factuelle (? ou token factuel)
    if is_factual_question(user_message):
        return True
    # Sinon : pas de RAG (navigation de menu)
    return False
\end{lstlisting}

\subsection{Ingestion idempotente}

L'ingestion est rendue idempotente par vérification préalable : si un document avec
le même \texttt{source\_uri} existe déjà avec \texttt{status='ready'}, il est ignoré
(sauf si \texttt{force=True}). En mode force, les anciens chunks sont supprimés de
PostgreSQL \textit{et} de Qdrant avant réingestion.

\begin{lstlisting}[language=Python, caption={Vérification d'idempotence dans ingest\_sources()},
                   label={lst:idempotent}]
existing = conn.execute(
    """SELECT id FROM kb_documents
       WHERE source_uri = %s AND status = 'ready'
       ORDER BY created_at DESC LIMIT 1""",
    (source_uri,),
).fetchone()

if existing and not force:
    stats["skipped"] += 1
    continue  # Document déjà indexé, on passe au suivant

if existing and force:
    # Purger l'ancienne version (DB + Qdrant)
    old_chunk_ids = _delete_document_chunks(conn, existing[0])
    delete_qdrant_points(old_chunk_ids)
    conn.execute("DELETE FROM kb_documents WHERE id = %s", (existing[0],))
    stats["updated"] += 1
\end{lstlisting}

%------------------------------------------------------------
\section{Sprint 3 — Sécurité et qualité logicielle}
%------------------------------------------------------------

\subsection{Authentification JWT HS256}

\subsubsection{Problème initial}

L'interface d'administration était protégée uniquement par un mot de passe statique
transmis en clair dans le header \texttt{X-Admin-Password}. Ce mécanisme ne permettait
ni expiration des sessions, ni audit des accès, et exposait le mot de passe en clair
dans les logs réseau.

\subsubsection{Implémentation JWT stdlib}

Afin d'éviter toute dépendance externe, le JWT HS256 a été implémenté avec la bibliothèque
standard Python (\texttt{hmac}, \texttt{hashlib}, \texttt{base64}).

\begin{lstlisting}[language=Python, caption={Implémentation JWT HS256 (stdlib Python)},
                   label={lst:jwt}]
def _encode_jwt() -> str:
    secret = os.getenv("JWT_SECRET") or os.getenv("ADMIN_PASSWORD")
    ttl    = int(os.getenv("JWT_TTL_MINUTES", "60"))
    header  = _b64url_encode(json.dumps({"alg": "HS256", "typ": "JWT"}).encode())
    now     = int(time.time())
    payload = _b64url_encode(json.dumps({
        "sub": "admin", "iat": now, "exp": now + ttl * 60
    }).encode())
    signing_input = f"{header}.{payload}"
    sig = _b64url_encode(
        hmac.new(secret.encode(), signing_input.encode(),
                 hashlib.sha256).digest()
    )
    return f"{signing_input}.{sig}"

def require_admin_password(
    x_admin_password: str | None = Header(None),
    authorization:    str | None = Header(None),
) -> None:
    # Mode JWT (recommandé)
    if authorization and authorization.startswith("Bearer "):
        _decode_jwt(authorization.removeprefix("Bearer ").strip())
        return
    # Rétrocompatibilité : mot de passe statique
    _ensure_admin_password(x_admin_password)
\end{lstlisting}

\paragraph{Variables d'environnement.}
\begin{itemize}
  \item \texttt{JWT\_SECRET} : secret HMAC (défaut = \texttt{ADMIN\_PASSWORD})
  \item \texttt{JWT\_TTL\_MINUTES} : durée de validité du token (défaut = 60 min)
  \item \texttt{ADMIN\_PASSWORD} : mot de passe de connexion admin
\end{itemize}

\subsection{Audit log}

Les actions sensibles génèrent des entrées structurées dans les logs applicatifs :

\begin{lstlisting}[language=Python, caption={Exemples d'entrées d'audit log},
                   label={lst:audit}]
# Connexion réussie
LOGGER.info("audit action=login status=success")

# Suppression d'un document KB
LOGGER.info("audit action=delete_kb_document document_id=%s title=%s chunks=%s",
            doc_id, title, len(chunk_ids))

# Lancement d'une ingestion
LOGGER.info("audit action=ingestion_start title=%s source_uri=%s source_type=%s",
            title, source_uri, source_type)
\end{lstlisting}

\subsection{Tests de non-régression}

Une suite de 24 tests unitaires a été développée dans \texttt{test\_sprint2.py},
couvrant les quatre fonctions modifiées au Sprint 2.

\begin{table}[H]
\centering
\caption{Résultats de la suite de tests non-régression (Sprint 2)}
\small
\begin{tabularx}{\textwidth}{|l|c|X|}
\hline
\rowcolor{tableheader}
\textbf{Module testé} & \textbf{Cas de test} & \textbf{Couverture} \\
\hline
\texttt{classify\_intent()} & 8 & Normalisation accents, best-score, regex bonus, absence de match \\
\hline
\texttt{chunk\_text()} & 8 & Paragraphes, overlap, texte long, vide, frontières préservées \\
\hline
\texttt{should\_trigger\_rag()} & 6 & Step wizard (False), intent RAG (True), question factuelle (True), navigation (False) \\
\hline
\texttt{ingest\_sources()} & 2 & Skip idempotent, force réingestion \\
\hline
\textbf{Total} & \textbf{24} & \textbf{Tous passing} \\
\hline
\end{tabularx}
\label{tab:tests}
\end{table}

Pour exécuter la suite de tests :

\begin{lstlisting}[language=bash, caption={Exécution des tests},
                   label={lst:pytest}]
# Depuis la racine du projet
cd backend
pytest tests/test_sprint2.py -v

# Résultat attendu
tests/test_sprint2.py::TestClassifyIntent::test_normalise_accent PASSED
tests/test_sprint2.py::TestChunkText::test_paragraph_boundary_preserved PASSED
tests/test_sprint2.py::TestShouldTriggerRag::test_false_on_form_step PASSED
# ... 24 tests PASSED in 0.8s
\end{lstlisting}

%------------------------------------------------------------
\section{Interface utilisateur}
%------------------------------------------------------------

\subsection{Interface de chat}

L'interface de chat Next.js offre une expérience conversationnelle guidée accessible
depuis un navigateur web standard.

Les fonctionnalités principales sont :
\begin{itemize}
  \item \textbf{Menu d'accueil} avec boutons d'orientation (Audience, Solutions, Budget,
  Rappel) permettant à l'utilisateur de choisir son parcours sans avoir à formuler une requête.
  \item \textbf{Parcours budgétaire} guidé : une série de questions (type d'annonceur,
  objectif marketing, tranche de budget) aboutit à une recommandation personnalisée.
  \item \textbf{Réponses RAG} pour les questions factuelles : le chatbot récupère les
  passages pertinents de la base documentaire et génère une réponse contextualisée.
  \item \textbf{Formulaires de leads} structurés : collecte des informations de contact
  avec validation des champs obligatoires.
  \item \textbf{Streaming} : les réponses du LLM s'affichent progressivement, améliorant
  la perception de réactivité.
\end{itemize}

\subsection{Comportement conversationnel — exemples}

\begin{table}[H]
\centering
\caption{Exemples de parcours conversationnels et comportements attendus}
\small
\begin{tabularx}{\textwidth}{|X|l|X|}
\hline
\rowcolor{tableheader}
\textbf{Message utilisateur} & \textbf{Step} & \textbf{Comportement} \\
\hline
« Bonjour » & WELCOME\_SCOPE & Présentation + boutons d'orientation \\
\hline
« Quel est le CPM du display ? » & MAIN\_MENU & RAG → réponse avec tarifs depuis KB \\
\hline
[Clic bouton Budget] & MAIN\_MENU & Transition vers BUDGET\_CLIENT\_TYPE \\
\hline
« Grande entreprise » & BUDGET\_CLIENT\_TYPE & Transition wizard suivante (pas de RAG) \\
\hline
« Je veux être rappelé » & n.a. & Déclenchement formulaire FORM\_STANDARD \\
\hline
\end{tabularx}
\label{tab:uc}
\end{table}

%------------------------------------------------------------
\section{Interface d'administration}
%------------------------------------------------------------

L'interface d'administration est accessible à l'URL \texttt{/admin}. Elle comprend
cinq onglets fonctionnels.

\subsection{Tableau de bord (Overview)}

La page d'accueil admin affiche les compteurs globaux :
\begin{itemize}
  \item Nombre de sessions de conversation initiées
  \item Nombre de messages échangés
  \item Nombre de fiches de contact collectées (leads)
  \item Nombre de documents KB indexés et de chunks générés
\end{itemize}

\subsection{Gestion de la base de connaissances}

L'onglet KB permet de visualiser les documents ingérés (titre, source, statut, nombre
de chunks) et de les supprimer. La suppression d'un document déclenche automatiquement
la purge des chunks associés dans PostgreSQL et Qdrant.

\subsection{Ingestion manuelle}

L'onglet Ingestion propose un workflow en plusieurs étapes :
\begin{enumerate}
  \item \textbf{Upload} : importation d'un fichier (\texttt{.txt}, \texttt{.md},
  \texttt{.json}, \texttt{.jsonl}, \texttt{.pdf})
  \item \textbf{Prévisualisation} : affichage du découpage en chunks avant indexation
  \item \textbf{Exécution} : lancement de l'ingestion complète avec logs temps réel (NDJSON)
\end{enumerate}

\subsection{Authentification et sécurité}

Depuis le Sprint 3, la connexion admin génère un token JWT retourné dans la réponse
de login. Toutes les requêtes suivantes utilisent ce token via le header
\texttt{Authorization: Bearer <token>}. Le token expire après 60 minutes (configurable
via \texttt{JWT\_TTL\_MINUTES}). La rétrocompatibilité avec \texttt{X-Admin-Password}
est maintenue pour la transition.

%------------------------------------------------------------
\section{Déploiement en production}
%------------------------------------------------------------

\subsection{Variables d'environnement clés}

\begin{table}[H]
\centering
\caption{Variables d'environnement principales}
\small
\begin{tabularx}{\textwidth}{|l|l|X|}
\hline
\rowcolor{tableheader}
\textbf{Variable} & \textbf{Valeur par défaut} & \textbf{Rôle} \\
\hline
\texttt{ADMIN\_PASSWORD} & \texttt{change\_me} & Mot de passe admin (à changer en prod) \\
\hline
\texttt{JWT\_SECRET} & = ADMIN\_PASSWORD & Secret HMAC pour signature JWT \\
\hline
\texttt{JWT\_TTL\_MINUTES} & \texttt{60} & Durée de validité du token JWT \\
\hline
\texttt{EMBEDDING\_URL} & localhost:8001 & Endpoint API embeddings Ollama \\
\hline
\texttt{RAG\_CHUNK\_SIZE} & \texttt{200} & Taille max d'un chunk (en mots) \\
\hline
\texttt{RAG\_LEXICAL\_WEIGHT} & \texttt{0.35} & Poids BM25 dans la fusion hybride \\
\hline
\end{tabularx}
\label{tab:env}
\end{table}

\subsection{Procédure de mise à jour en production}

\begin{lstlisting}[language=bash, caption={Procédure de mise à jour sur le serveur},
                   label={lst:deploy}]
# Récupération des dernières modifications
git pull origin main

# Rebuild et redémarrage du backend et frontend
docker compose -f infra/docker-compose.yml build backend frontend
docker compose -f infra/docker-compose.yml up -d backend frontend

# Si modification du schéma de données (migrations SQL)
docker compose -f infra/docker-compose.yml exec backend \
    python -m scripts.migrate

# Réingestion KB après modification des fichiers sources
docker compose -f infra/docker-compose.yml exec backend \
    python -c "from app.rag.ingest import ingest_sources; \
               ingest_sources(force=True)"
\end{lstlisting}

%------------------------------------------------------------
\section{Conclusion}
%------------------------------------------------------------

Les trois sprints ont permis de livrer progressivement un système conversationnel
complet et opérationnel. Le Sprint 1 a posé l'infrastructure technique. Le Sprint 2 a
significativement amélioré la qualité de l'IA en rendant le classifieur plus robuste,
le chunking plus cohérent et le pipeline RAG plus efficace. Le Sprint 3 a sécurisé
l'accès administratif et garanti la non-régression via une suite de tests automatisés.
L'ensemble forme une solution déployable, maintenable et évolutive.

%============================================================
% CONCLUSION ET PERSPECTIVES
%============================================================
\chapter*{Conclusion et perspectives}
\addcontentsline{toc}{chapter}{Conclusion et perspectives}
\markboth{Conclusion et perspectives}{}

\lettrine[lines=2]{C}{e} projet a permis de concevoir et de déployer \textbf{TNChatbot},
un assistant conversationnel intelligent pour Tunisie Numérique. En combinant un parcours
guidé par machine à états, un moteur RAG hybride (vectoriel + BM25) et un modèle de
langage hébergé localement via Ollama, la solution répond concrètement aux besoins
identifiés : qualifier les prospects, présenter les offres publicitaires et collecter des leads
exploitables pour l'équipe commerciale.

\paragraph{Bilan technique.}
Les cinq cas d'usage principaux (découverte de l'audience, exploration des solutions,
recommandation budget, questions factuelles RAG, capture de leads) sont opérationnels
et testés. L'architecture modulaire microservices garantit la maintenabilité et facilite
l'évolution du système. La suite de 24 tests non-régression protège les fonctionnalités
clés contre les régressions futures.

\paragraph{Limites identifiées.}
La qualité des réponses RAG reste directement liée à la richesse et à la mise à jour
régulière des fichiers sources dans \texttt{kb\_sources/}. Par ailleurs, l'utilisation d'un
LLM local impose des contraintes de ressources matérielles qui peuvent impacter les
performances en charge.

\paragraph{Perspectives d'évolution.}
Plusieurs axes d'amélioration sont envisageables à court et moyen terme :
\begin{itemize}
  \item \textbf{Observabilité} : intégration d'un tableau de bord Prometheus/Grafana pour
  surveiller les métriques de performance en temps réel (latence, taux de fallback,
  taux de conversion en lead).
  \item \textbf{RBAC admin} : ajout d'un système de rôles (éditeur, opérateur, administrateur)
  sur les endpoints sensibles de l'interface d'administration.
  \item \textbf{Reranker cross-encoder} : activation du reranking Cohere pour améliorer
  la précision du retrieval RAG sur les questions complexes.
  \item \textbf{Évaluation offline du RAG} : mise en place d'un jeu de questions/réponses
  de référence pour mesurer Recall@k et MRR avant et après chaque modification du
  pipeline de retrieval.
  \item \textbf{Support multilingue} : extension du classifieur d'intention et de la base
  documentaire pour gérer les requêtes en arabe, fréquentes sur le marché tunisien.
\end{itemize}

En somme, ce projet pose les bases d'une solution intelligente, évolutive et opérationnelle
pour la gestion automatisée de la relation commerciale sur Tunisie Numérique, et constitue
un exemple concret d'intégration des technologies RAG dans un contexte métier réel.

%============================================================
% BIBLIOGRAPHIE
%============================================================
\chapter*{Bibliographie}
\addcontentsline{toc}{chapter}{Bibliographie}
\markboth{Bibliographie}{}

\begin{enumerate}[label={[\arabic*]}, leftmargin=2.5em]
  \item Lewis, P. et al. \textit{Retrieval-Augmented Generation for
  Knowledge-Intensive NLP Tasks}. Advances in Neural Information Processing Systems
  (NeurIPS), 2020.

  \item Robertson, S., \& Zaragoza, H. \textit{The Probabilistic Relevance Framework:
  BM25 and Beyond}. Foundations and Trends in Information Retrieval, 3(4), 333--389, 2009.

  \item Cormack, G. V., Clarke, C. L. A., \& Buettcher, S. \textit{Reciprocal Rank
  Fusion Outperforms Condorcet and Individual Rank Learning Methods}. SIGIR, 2009.

  \item FastAPI Documentation. \url{https://fastapi.tiangolo.com}. (Dernier accès :
  Juin 2025).

  \item Next.js Documentation. \url{https://nextjs.org/docs}. (Dernier accès :
  Juin 2025).

  \item Qdrant Documentation. \url{https://qdrant.tech/documentation}. (Dernier accès :
  Juin 2025).

  \item Ollama GitHub Repository. \url{https://github.com/ollama/ollama}. (Dernier
  accès : Juin 2025).

  \item Docker Compose Documentation.
  \url{https://docs.docker.com/compose}. (Dernier accès : Juin 2025).

  \item PostgreSQL Documentation. \url{https://www.postgresql.org/docs}. (Dernier
  accès : Juin 2025).

  \item RFC 7519 — JSON Web Token (JWT). \url{https://datatracker.ietf.org/doc/html/rfc7519}.
  (Dernier accès : Juin 2025).
\end{enumerate}

\end{document}
